\documentclass[11pt,a4paper]{article}
\usepackage{gtu-el-yazisi}
\begin{document}
\notbilgi{İşaretler ve Sistemler}{Ders 1 -- İşaretler: enerji, güç ve temel işlemler}{Isaret-ve-Sistemler-ders-1.pdf}

\section{İşaretler ve Sistemler}
\sayfa{1}
\begin{tanim}[(İşaret -- signal)]
  Herhangi bir fiziki olay hakkında bilgi taşıyan, bir veya daha fazla bağımsız değişkeni olan fonksiyondur.
\end{tanim}

\subsection{Sürekli zamanlı işaret (continuous-time signal)}
\begin{minipage}[c]{0.6\linewidth}
  $x(t)$\\
  \quad $\hookrightarrow$ bağımsız değişkeni $t$ süreklidir.
\end{minipage}\hfill
\begin{minipage}[c]{0.36\linewidth}
  \centering
  \begin{tikzpicture}[scale=0.8]
    \draw[->] (-1.8,0) -- (2,0) node[right] {$t$};
    \draw[->] (0,-0.8) -- (0,1.2) node[above] {$x(t)$};
    \draw[murekkep, thick, smooth] plot[domain=-1.6:1.6, samples=60] (\x, {0.6*sin(deg(3*\x))+0.25*sin(deg(5*\x+1))});
  \end{tikzpicture}
\end{minipage}

ör: Ses işareti akustik basıncın zamana bağlı bir fonksiyonudur.\\
\phantom{ör: }Resim parlaklığın iki uzamsal $(x,y)$ değişkene bağlı bir fonksiyonudur.

\subsection{Ayrık zamanlı işaret (discrete-time signal)}
\begin{minipage}[c]{0.6\linewidth}
  $x[n]$\\
  \quad $\hookrightarrow$ bağımsız değişkeni $n$ süreksizdir ($n\in\mathbb{Z}$, tamsayı değerleri alır.)
\end{minipage}\hfill
\begin{minipage}[c]{0.36\linewidth}
  \centering
  \begin{tikzpicture}[scale=0.75]
    \draw[->] (-2.8,0) -- (2.8,0) node[right] {$n$};
    \draw[->] (0,-0.7) -- (0,1.7) node[above] {$x[n]$};
    \foreach \k in {-2,-1,1,2} \draw (\k,0.06) -- (\k,-0.06) node[below] {\small$\k$};
    \node[below left] at (0,0) {\small$0$};
    \draw[murekkep, thick] (-2,0) -- (-2,1.2);
    \fill[murekkep] (-2,1.2) circle (2pt) node[above] {$x[-2]$};
    \draw[murekkep, thick] (1,0) -- (1,0.8);
    \fill[murekkep] (1,0.8) circle (2pt) node[right] {$x[1]$};
    \draw[dashed, kursunkalem] (0,0.8) -- (1,0.8);
  \end{tikzpicture}
\end{minipage}

ör: Yıllık nüfus sayısı

Ayrık zamanlı işaret;
\begin{itemize}[label=--]
  \item olayın doğası gereği ayrık değerler alabilir.
  \item sürekli zamanlı işaretin zamanda belli aralıklarla örneklenmesiyle elde edilebilir.
\end{itemize}
\[ x(t) \xrightarrow[\text{Sampling}]{@\,nT_s} x(nT_s) \equiv x[n] \]

\subsection{Enerji (energy of a signal)}
\begin{align*}
  E_x &= \int_{-\infty}^{\infty} |x(t)|^2\,dt = \lim_{T\to\infty} \int_{-T}^{T} |x(t)|^2\,dt \\
  E_x &= \sum_{n=-\infty}^{\infty} |x[n]|^2 = \lim_{N\to\infty} \sum_{n=-N}^{N} |x[n]|^2
\end{align*}

\subsection{Güç (power of a signal)}
\begin{align*}
  P_x &= \lim_{T\to\infty} \frac{1}{2T} \int_{-T}^{T} |x(t)|^2\,dt \\
  P_x &= \lim_{N\to\infty} \frac{1}{2N+1} \sum_{n=-N}^{N} |x[n]|^2
\end{align*}

\begin{itemize}[label=$*$]
  \item $E_x<\infty$ ise $x(t)/x[n]$ işareti enerji işaretidir ve ortalama gücü sıfırdır, $P_x=0$.
  \item \vurgu[turuncu]{$0<P_x<\infty$} ise $x(t)/x[n]$ işareti güç işaretidir ve enerjisi sonsuzdur, $E_x\to\infty$.
  \item Bazı işaretler ne enerji ne de güç işareti olmayabilirler.\\
  Ör: $x(t)=kt \;\longrightarrow\; E_x\to\infty$\\
  \phantom{Ör: $x(t)=kt$\;}$\hookrightarrow P_x\to\infty$
\end{itemize}

\section{Bağımsız değişkenlerin dönüşümleri}
\sayfa{2}
(Transformation of independent variable)

\subsection{Zamanda öteleme (time shifting)}
\[
  x(t)\longrightarrow x(t-t_0)\;
  \begin{array}{l}
    \longrightarrow t_0>0 \ \ \text{gecikme (delay), sağa öteleme}\\
    \hookrightarrow t_0<0 \ \ \text{ilerleme (advance), sola öteleme}
  \end{array}
\]
\[ x[n]\longrightarrow x[n-n_0] \]

Ex:\par
\begin{center}
\begin{tikzpicture}[scale=0.5]
  \draw[->] (-1.8,0) -- (2.9,0) node[right] {$t$};
  \draw[->] (0,-0.6) -- (0,2.9) node[above] {$x(t)$};
  \draw[murekkep, thick] (-1,0) -- (0,2) -- (2,2) -- (2,0);
  \foreach \k in {-1,1,2} \draw (\k,0.08) -- (\k,-0.08) node[below] {\scriptsize$\k$};
  \node[below left] at (0,0) {\scriptsize$0$};
  \node[left] at (0,2) {\scriptsize$2$};
\end{tikzpicture}\hfill
\begin{tikzpicture}[scale=0.5]
  \draw[->] (-1,0) -- (3.9,0) node[right] {$t$};
  \draw[->] (0,-0.6) -- (0,2.9) node[above] {$x(t-1)$};
  \draw[murekkep, thick] (0,0) -- (1,2) -- (3,2) -- (3,0);
  \foreach \k in {1,2,3} \draw (\k,0.08) -- (\k,-0.08) node[below] {\scriptsize$\k$};
  \node[left] at (0,2) {\scriptsize$2$};
  \node[below, align=center, font=\small] at (1.5,-1.1) {$t_0=1>0$ gecikme/sağa öteleme};
\end{tikzpicture}\hfill
\begin{tikzpicture}[scale=0.5]
  \draw[->] (-2.8,0) -- (1.9,0) node[right] {$t$};
  \draw[->] (0,-0.6) -- (0,2.9) node[above] {$x(t+1)$};
  \draw[murekkep, thick] (-2,0) -- (-1,2) -- (1,2) -- (1,0);
  \foreach \k in {-2,-1,1} \draw (\k,0.08) -- (\k,-0.08) node[below] {\scriptsize$\k$};
  \node[above right] at (0,2) {\scriptsize$2$};
  \node[below, align=center, font=\small] at (-0.4,-1.1) {$t_0=-1<0$ ilerleme/sola öteleme};
\end{tikzpicture}
\end{center}

\subsection{Zamanda ters çevirme (time reversal)}
\[ x(t)\longrightarrow x(-t)\longrightarrow \text{orijine veya sıfır noktasına göre yansımasını alma} \]
\[ x[n]\longrightarrow x[-n] \]

Ex:\par
\begin{center}
\begin{tikzpicture}[scale=0.5]
  \draw[->] (-1.8,0) -- (2.9,0) node[right] {$t$};
  \draw[->] (0,-0.6) -- (0,2.9) node[above] {$x(t)$};
  \draw[murekkep, thick] (-1,0) -- (0,2) -- (2,2) -- (2,0);
  \foreach \k in {-1,2} \draw (\k,0.08) -- (\k,-0.08) node[below] {\scriptsize$\k$};
  \node[below right] at (0,0) {\scriptsize$0$};
  \node[left] at (0,2) {\scriptsize$2$};
\end{tikzpicture}\qquad\qquad
\begin{tikzpicture}[scale=0.5]
  \draw[->] (-2.8,0) -- (1.9,0) node[right] {$t$};
  \draw[->] (0,-0.6) -- (0,2.9) node[above] {$x(-t)$};
  \draw[murekkep, thick] (-2,0) -- (-2,2) -- (0,2) -- (1,0);
  \foreach \k in {-2,-1,1} \draw (\k,0.08) -- (\k,-0.08) node[below] {\scriptsize$\k$};
  \node[above right] at (0,2) {\scriptsize$2$};
\end{tikzpicture}
\end{center}

\subsection{Zamanda ölçekleme (time scaling)}
\[
  x(t)\longrightarrow x(\alpha t)\;
  \begin{array}{l}
    \longrightarrow \alpha>1 \ \ \text{sıkıştırma (compression)}\\
    \hookrightarrow 0<\alpha<1 \ \ \text{genişletme (expansion)}
  \end{array}
\]
\[ x[n]\longrightarrow x[\alpha n] \]

Ex:\par
\noindent
\begin{minipage}[t]{0.28\linewidth}
  \centering
  \begin{tikzpicture}[scale=0.5]
    \draw[->] (-1.8,0) -- (2.9,0) node[right] {$t$};
    \draw[->] (0,-0.6) -- (0,2.9) node[above] {$x(t)$};
    \draw[murekkep, thick] (-1,0) -- (0,2) -- (2,2) -- (2,0);
    \foreach \k in {-1,1,2} \draw (\k,0.08) -- (\k,-0.08) node[below] {\scriptsize$\k$};
    \node[left] at (0,2) {\scriptsize$2$};
  \end{tikzpicture}
\end{minipage}\hfill
\begin{minipage}[t]{0.34\linewidth}
  \centering
  \begin{tikzpicture}[scale=0.5]
    \draw[->] (-1.8,0) -- (2.9,0) node[right] {$t$};
    \draw[->] (0,-0.6) -- (0,2.9) node[above] {$x(2t)$};
    \draw[murekkep, thick] (-1,0) -- (0,2) -- (2,2) -- (2,0);
    \draw (-1,0.08) -- (-1,-0.08) node[below] {\scriptsize$-\frac12$};
    \draw (1,0.08) -- (1,-0.08) node[below] {\scriptsize$\frac12$};
    \draw (2,0.08) -- (2,-0.08) node[below] {\scriptsize$1$};
    \node[left] at (0,2) {\scriptsize$2$};
  \end{tikzpicture}\\[1mm]
  \raggedright
  $x(t)$, $t:[-1,2]$\\
  $x(2t)$, $2t:[-1,2]$\\
  \quad$\hookrightarrow t:[-1/2,\,1]$\\
  $\alpha=2>1$ sıkıştırma
\end{minipage}\hfill
\begin{minipage}[t]{0.34\linewidth}
  \centering
  \begin{tikzpicture}[scale=0.5]
    \draw[->] (-2.8,0) -- (4.9,0) node[right] {$t$};
    \draw[->] (0,-0.6) -- (0,2.9) node[above] {$x(t/2)$};
    \draw[murekkep, thick] (-2,0) -- (0,2) -- (4,2) -- (4,0);
    \foreach \k in {-2,2,4} \draw (\k,0.08) -- (\k,-0.08) node[below] {\scriptsize$\k$};
    \node[left] at (0,2) {\scriptsize$2$};
  \end{tikzpicture}\\[1mm]
  \raggedright
  $x(t)$, $t:[-1,2]$\\
  $x(t/2)$, $t/2:[-1,2]$\\
  \quad$\hookrightarrow t:[-2,4]$\\
  $0<\alpha=\frac12<1$ genişletme
\end{minipage}

\sayfa{3}
\[
  *\ \ x(t)\longrightarrow x(\alpha t+\beta)\qquad
  \begin{array}{l}
    \text{1) } \beta \text{ kadar öteleme}\\
    \text{2) } \alpha \text{ kadar ölçekleme ve gerekiyorsa zamanda ters çevirme}
  \end{array}
\]

\begin{ornek}[]
  $x(t)$ veriliyor; \quad $x(-2t+5)=\,?$\par
  \begin{center}
  \begin{tikzpicture}[scale=0.5]
    \draw[->] (-1.8,0) -- (2.9,0) node[right] {$t$};
    \draw[->] (0,-0.6) -- (0,2.9) node[above] {$x(t)$};
    \draw[murekkep, thick] (-1,0) -- (0,2) -- (2,2) -- (2,0);
    \foreach \k in {-1,1,2} \draw (\k,0.08) -- (\k,-0.08) node[below] {\scriptsize$\k$};
    \node[below right] at (0,0) {\scriptsize$0$};
    \node[left] at (0,2) {\scriptsize$2$};
  \end{tikzpicture}
  \end{center}
  \begin{center}
  \begin{tikzpicture}[scale=0.45]
    \draw[->] (-6.8,0) -- (1,0) node[right] {$t$};
    \draw[->] (0,-0.6) -- (0,2.9) node[above] {$x(t+5)$};
    \draw[murekkep, thick] (-6,0) -- (-5,2) -- (-3,2) -- (-3,0);
    \foreach \k in {-6,-5,-4,-3} \draw (\k,0.08) -- (\k,-0.08) node[below] {\scriptsize$\k$};
    \draw (-0.08,2) -- (0.08,2) node[right] {\scriptsize$2$};
  \end{tikzpicture}\hfill
  \begin{tikzpicture}[scale=0.45]
    \draw[->] (-0.8,0) -- (6.9,0) node[right] {$t$};
    \draw[->] (0,-0.6) -- (0,2.9) node[above] {$x(-t+5)$};
    \draw[murekkep, thick] (3,0) -- (3,2) -- (5,2) -- (6,0);
    \foreach \k in {3,4,5,6} \draw (\k,0.08) -- (\k,-0.08) node[below] {\scriptsize$\k$};
    \draw (-0.08,2) -- (0.08,2); \node[left] at (0,2) {\scriptsize$2$};
  \end{tikzpicture}\hfill
  \begin{tikzpicture}[scale=0.45]
    \draw[->] (-0.8,0) -- (6.9,0) node[right] {$t$};
    \draw[->] (0,-0.6) -- (0,2.9) node[above] {$x(-2t+5)$};
    \draw[murekkep, thick] (3,0) -- (3,2) -- (5,2) -- (6,0);
    \draw[dashed, kursunkalem] (0,2) -- (3,2);
    \draw (3,0.08) -- (3,-0.08) node[below] {\scriptsize$\frac32$};
    \draw (4,0.08) -- (4,-0.08) node[below] {\scriptsize$2$};
    \draw (5,0.08) -- (5,-0.08) node[below] {\scriptsize$\frac52$};
    \draw (6,0.08) -- (6,-0.08) node[below] {\scriptsize$3$};
    \node[left] at (0,2) {\scriptsize$2$};
  \end{tikzpicture}
  \end{center}
  \begin{align*}
    &x(t),\quad t:[-1,2]\\
    &x(-2t+5),\quad -2t+5:[-1,2]\\
    &\phantom{x(-2t+5),\quad -2t+5:\ }t:[1{,}5,\ 3]
  \end{align*}
\end{ornek}

\begin{ornek}[]
  $x[n]=\begin{cases} 0, & n<-2 \text{ ve } n>3\\ 1, & \text{diğer}\end{cases}$
  \ olmak üzere \ $x[-n/2+1]=\,?$
\end{ornek}
\begin{center}
  \begin{tikzpicture}[scale=0.42]
    \draw[->] (-3,0) -- (4,0) node[right] {$n$};
    \draw[->] (0,-0.4) -- (0,2.2) node[above] {$x[n]$};
    \draw[dashed, kursunkalem] (-2,1.2) -- (3,1.2);
    \foreach \k in {-2,...,3} {\draw[murekkep, thick] (\k,0) -- (\k,1.2); \fill[murekkep] (\k,1.2) circle (3pt);}
    \foreach \k in {-2,-1,0,1,2,3} \node[below] at (\k,0) {\scriptsize$\k$};
    \node[above left] at (0,1.2) {\scriptsize$1$};
  \end{tikzpicture}\hspace{1.5cm}
  \begin{tikzpicture}[scale=0.42]
    \draw[->] (-5,0) -- (7,0) node[right] {$n$};
    \draw[->] (0,-0.4) -- (0,2.2) node[above] {$x[-n/2+1]$};
    \draw[dashed, kursunkalem] (-4,1.2) -- (6,1.2);
    \foreach \k in {-4,...,6} {\draw[murekkep, thick] (\k,0) -- (\k,1.2); \fill[murekkep] (\k,1.2) circle (3pt);}
    \foreach \k in {-4,...,6} \node[below] at (\k,0) {\tiny$\k$};
  \end{tikzpicture}
\end{center}
$x[n],\ n:[-2,3]$\\
$x[-n/2+1],\ -\frac{n}{2}+1:[-2,3]$\\
\phantom{$x[-n/2+1],\ -\frac{n}{2}+1:$\,}$n:[-4,6]$\footnote{Claude notu: ayrık zamanda $-n/2+1$ yalnızca çift $n$ için tamsayıdır; tek $n$ değerlerinde $x[-n/2+1]$ tanımsızdır (genellikle $0$ alınır), yani çubuklar muhtemelen yalnızca $n=-4,-2,0,2,4,6$ noktalarında olmalı.}

\section{Periyodik işaretler (periodic signals)}
\sayfa{4}
\begin{align*}
  x(t) &= x(t+kT_0), && t\in\mathbb{R},\ k\in\mathbb{Z},\ T_0:\text{ periyot} && \longrightarrow \text{Periyodik.}\\
  x[n] &= x[n+kN_0], && n,k\in\mathbb{Z},\ N_0:\text{ periyot}
\end{align*}
Bu eşitlikleri sağlayan $T_0/N_0$ yok ise işaret aperiyodiktir.

\noindent
\begin{minipage}[c]{0.55\linewidth}
  Ex:\par
  \centering
  \begin{tikzpicture}[scale=0.7]
    \draw[->] (-4,0) -- (4,0) node[right] {$t$};
    \draw[->] (0,-0.3) -- (0,1.7) node[above] {$x(t)$};
    \foreach \c in {-2.6,0,2.6} \draw[murekkep, thick] (\c-0.9,0) -- (\c,1.2) -- (\c+0.9,0);
    \draw[dashed, kursunkalem] (-2.6,0) -- (-2.6,1.2);
    \draw[dashed, kursunkalem] (2.6,0) -- (2.6,1.2);
    \node[below] at (-2.6,0) {\small$-T$};
    \node[below] at (0,0) {\small$0$};
    \node[below] at (2.6,0) {\small$T$};
    \node at (-3.7,0.5) {$\cdots$};
    \node at (3.7,0.5) {$\cdots$};
  \end{tikzpicture}
\end{minipage}\hfill
\begin{minipage}[c]{0.42\linewidth}
  $\Rightarrow x(t)=x(t+T)$, \ $T$ ile periyodiktir.
\end{minipage}

\medskip
Ex: $x(t)=\cos(t)$
\[
  x(t+kT_0)=\cos(t+kT_0)\underset{T_0=2\pi}{=}\cos(t+2\pi k)=\cos(t)=x(t)
\]

\section{Tek (Odd) / Çift (Even) İşaretler}
\noindent
\begin{minipage}[c]{0.66\linewidth}
  \underline{Çift:} \ $x(t)=x(-t)\longrightarrow \forall t$ için $y$ eksenine göre simetrik ise çift işarettir.\\
  \phantom{\underline{Çift:} \ }$x[n]=x[-n]$
\end{minipage}\hfill
\begin{minipage}[c]{0.3\linewidth}
  \centering
  \begin{tikzpicture}[scale=0.55]
    \draw[->] (-2.6,0) -- (2.8,0) node[right] {$t$};
    \draw[->] (0,-0.3) -- (0,1.8) node[above] {$x(t)$};
    \draw[murekkep, thick, smooth] plot[domain=-2:2, samples=61] (\x, {abs(sin(deg(pi*\x/2)))});
    \foreach \k in {-2,-1,1,2} \draw (\k,0.06) -- (\k,-0.06) node[below] {\scriptsize$\k$};
    \node[below left] at (0,0) {\scriptsize$0$};
    \node[left] at (0,1) {\scriptsize$1$};
  \end{tikzpicture}
\end{minipage}

\noindent
\begin{minipage}[c]{0.66\linewidth}
  \underline{Tek:} \ $x(t)=-x(-t)\longrightarrow \forall t$ için orijine göre simetrik ise tek işarettir.\\
  \phantom{\underline{Tek:} \ }$x[n]=-x[-n]$
\end{minipage}\hfill
\begin{minipage}[c]{0.3\linewidth}
  \centering
  \begin{tikzpicture}[scale=0.55]
    \draw[->] (-2.6,0) -- (2.8,0) node[right] {$t$};
    \draw[->] (0,-1.3) -- (0,1.6) node[above] {$x(t)$};
    \draw[murekkep, thick, smooth] plot[domain=-2:2, samples=61] (\x, {sin(deg(pi*\x/2))});
    \draw[dashed, kursunkalem] (1,0) -- (1,1);
    \draw[dashed, kursunkalem] (-1,0) -- (-1,-1);
    \foreach \k in {-2,-1,1,2} \draw (\k,0.06) -- (\k,-0.06);
    \foreach \k in {-2,-1} \node[above] at (\k,0) {\scriptsize$\k$};
    \foreach \k in {1,2} \node[below] at (\k,0) {\scriptsize$\k$};
  \end{tikzpicture}
\end{minipage}

\[ *\ \ x(t)=x_e(t)+x_o(t) \]
\[
  x_e(t)=\frac{x(t)+x(-t)}{2}\qquad\qquad x_o(t)=\frac{x(t)-x(-t)}{2}
\]

\section{Üstel ve Sinüsoidal İşaretler (Exponential and sinusoidal Signals)}
\sayfa{5}
\subsection{Sürekli zamanlı sinüsoidal işaretler}
\[
  x(t)=A\cos(\omega_0 t+\phi)\longrightarrow
  \begin{array}{l}
    A:\ \text{genlik}\\
    \omega_0:\ \text{açısal frekans, [rad/sn]}\\
    \phi:\ \text{faz}
  \end{array}
\]
\[
  \omega_0=2\pi f_0=\frac{2\pi}{T_0}\ \longrightarrow\
  \begin{array}{l}
    T_0:\ \text{periyot, [sn]}\\[1mm]
    f_0=\dfrac{1}{T_0}:\ \text{frekans, } \left[\text{Hertz}=\dfrac{1}{\text{sn}}\right]
  \end{array}
\]

\noindent
\begin{minipage}[c]{0.42\linewidth}
  \centering
  \begin{tikzpicture}[scale=0.5]
    \draw[->] (-1.8,0) -- (8.2,0) node[right] {$t$};
    \draw[->] (0,-1.8) -- (0,2.2) node[above] {$A\cos(\omega_0t+\phi)$};
    \draw[murekkep, thick, smooth] plot[domain=-1.6:7.6, samples=90] (\x, {1.3*cos(deg(\x-0.6))});
    \draw[dashed, kursunkalem] (-0.1,1.3) -- (0.6,1.3);
    \node[left] at (0,1.3) {\scriptsize$A$};
    \node[left] at (0,0.75) {\scriptsize$A\cos\phi$};
    \draw[<->] (0.6,1.5) -- (6.88,1.5) node[midway, above] {\scriptsize$T_0$};
  \end{tikzpicture}
\end{minipage}\hfill
\begin{minipage}[c]{0.55\linewidth}
  $\Rightarrow$ (Açısal) frekans $\times$ Periyot\\
  \quad $\omega_0\!\uparrow$ \ $f_0\!\uparrow$ \ $T_0\!\downarrow$ \ işaretin salınımı hızlanır.\\[1mm]
  {\color{kursunkalem}\small $\omega_0\!\uparrow$ \ $f_0\!\uparrow$ \ $T_0\!\downarrow$ \ işaretin osilasyonu artar, periyot azalır\\
  $x(t)=\text{Constant}$ ise \ $T_0\to\infty$ (undefined) \ $f_0\to0$ osilasyonu yoktur}
\end{minipage}

\subsection{Sürekli zamanlı karmaşık üstel işaretler}
\[
  x(t)=Ce^{at}\longrightarrow
  \begin{array}{l}
    C=|C|e^{j\theta}\\
    a=r+j\omega_0
  \end{array}
\]
$*$ $a$ ve $C$ reel/gerçek ise $x(t)$ reel üstel (real exponential) işarettir.\\
\phantom{$*$ }$C=|C|$ ve $a=r$

\begin{center}
\begin{tikzpicture}[scale=0.6]
  \draw[->] (-3,0) -- (3.2,0) node[right] {$t$};
  \draw[->] (0,-0.2) -- (0,3.2) node[above] {$x(t)=Ce^{at}$};
  \draw[murekkep, thick, smooth] plot[domain=-3:2.6, samples=60] (\x, {0.6*exp(0.62*\x)});
  \node[left] at (0,0.6) {\scriptsize$C$};
  \node[below] at (0,-0.5) {\small$a>0$ iken artan üstel};
\end{tikzpicture}\qquad\qquad
\begin{tikzpicture}[scale=0.6]
  \draw[->] (-3,0) -- (3.2,0) node[right] {$t$};
  \draw[->] (0,-0.2) -- (0,3.2) node[above] {$x(t)$};
  \draw[murekkep, thick, smooth] plot[domain=-2.6:3, samples=60] (\x, {0.6*exp(-0.62*\x)});
  \node[right] at (0,0.6) {\scriptsize$C$};
  \node[below] at (0,-0.5) {\small$a<0$ iken azalan üstel};
\end{tikzpicture}
\end{center}

$*$ $a=j\omega_0$ $(r=0)$ olmak üzere tamamen sanal olduğunda $x(t)=e^{j\omega_0t}$ olup, periyodik karmaşık üstel (periodic complex exponential) işarettir.

\medskip
$x(t)\overset{?}{=}x(t+T)$ ise periyodiktir.\\
İspat: \ $x(t)=e^{j\omega_0t}$
\begin{align*}
  x(t+T)&=e^{j\omega_0(t+T)}=e^{j\omega_0t}\cdot e^{j\omega_0T}
\end{align*}
Eğer $e^{j\omega_0T}=1$ olursa $x(t)$'nin periyodikliği ispatlanır.
\[
  e^{j\omega_0T}=1=e^{j2\pi},\quad \omega_0T=2\pi\ \longrightarrow\ T=\frac{2\pi}{\omega_0},\ \ \omega_0\neq0
  \qquad
  {\color{kirmizikalem}\fbox{\begin{tabular}{l}$\omega_0=0$, $x(t)=1$\\ $T\to\infty$\end{tabular}}}
\]

\sayfa{6}
\noindent
\begin{minipage}[c]{0.72\linewidth}
  Euler bağıntısı $\Rightarrow e^{j\theta}=\cos\theta+j\sin\theta$
  \[ e^{j2\pi}=\underbrace{\cos2\pi}_{1}+j\underbrace{\sin2\pi}_{0}=1 \]
\end{minipage}\hfill
\begin{minipage}[c]{0.25\linewidth}
  \centering
  {\color{kirmizikalem}\fbox{$e^{j\pi}=-1$}}
\end{minipage}
\[
  e^{j\omega_0t}=\cos\omega_0t+j\sin\omega_0t\ \
  \begin{array}{l}
    \longrightarrow \cos\omega_0t=\bigl(e^{j\omega_0t}+e^{-j\omega_0t}\bigr)/2\\[1mm]
    \hookrightarrow \sin\omega_0t=\bigl(e^{j\omega_0t}-e^{-j\omega_0t}\bigr)/2j
  \end{array}
\]
\begin{align*}
  Ae^{j(\omega_0t+\phi)}&=A\cos(\omega_0t+\phi)+jA\sin(\omega_0t+\phi)\quad\longrightarrow\\
  A\cos(\omega_0t+\phi)&=\mathrm{Re}\bigl\{Ae^{j(\omega_0t+\phi)}\bigr\}\\
  A\sin(\omega_0t+\phi)&=\mathrm{Im}\bigl\{Ae^{j(\omega_0t+\phi)}\bigr\}
\end{align*}

\subsection{Ayrık zamanlı sinüsoidal işaretler}
\[
  x[n]=A\cos(\omega_0n+\phi)=\frac{A}{2}\Bigl[e^{j(\omega_0n+\phi)}+e^{-j(\omega_0n+\phi)}\Bigr]
  =\frac{A}{2}\Bigl[e^{j\phi}e^{j\omega_0n}+e^{-j\phi}e^{-j\omega_0n}\Bigr]
\]

\subsection{Ayrık zamanlı karmaşık üstel işaretler}
\noindent
\begin{minipage}[c]{0.65\linewidth}
  \[
    x[n]=C\alpha^n\longrightarrow
    \begin{array}{l}
      C=|C|e^{j\theta}\\
      \alpha=|\alpha|e^{j\omega_0}
    \end{array}
  \]
\end{minipage}\hfill
\begin{minipage}[c]{0.3\linewidth}
  \centering
  {\color{kirmizikalem}\fbox{$N_0=2\pi/\omega_0$}}
\end{minipage}

$*$ $C$ ve $\alpha$ reel ise $x[n]$ reel üstel işarettir.\\
$*$ $\alpha=e^{j\omega_0}$ olmak üzere $x[n]=e^{j\omega_0n}$ işareti periyodik üstel işarettir.

\subsection{Ayrık zamanlı karmaşık üstel işaretlerin periyodikliği}
$*$ $x[n]=e^{j\omega_0n}$
\[
  \omega_0+2\pi\longrightarrow e^{j(\omega_0+2\pi)n}=e^{j\omega_0n}\cdot\underbrace{e^{j2\pi n}}_{\to1}=e^{j\omega_0n}
\]
Ayrık zamanlı karmaşık üstel işaretler ancak $2\pi$ aralığındaki frekanslarda farklılık gösterirler.
Bu nedenle bu işaretleri $-\pi<\omega_0<\pi$ veya $0<\omega_0<2\pi$ aralıklarında incelemek yeterlidir.

\medskip
$*$ $x[n]\overset{?}{=}x[n+N]$
\begin{align*}
  x[n+N]&=e^{j\omega_0(n+N)}=e^{j\omega_0n}\cdot e^{j\omega_0N}\\
  e^{j\omega_0N}&=1=e^{j2\pi m}\ \longrightarrow\ \omega_0N=2\pi m\\
  &\phantom{=1=e^{j2\pi m}\ \longrightarrow\ }\frac{\omega_0}{2\pi}=\frac{m}{N},\ \ m \text{ ve } N>0\in\mathbb{Z}
    \ \Bigl(\tfrac{m}{N} \text{ kesirli yani iki tam sayının oranı olmalı}\Bigr)\\
  &\phantom{=1=e^{j2\pi m}\ \longrightarrow\ }N=\Bigl(\frac{2\pi}{\omega_0}\Bigr)m
\end{align*}

\sayfa{7}
\noindent
\begin{minipage}[t]{0.44\linewidth}
  \underline{$e^{j\omega_0t}$}\par\smallskip
  $\omega_0$'ın farklı değerleri için farklı işaret\\[1mm]
  Herhangi bir $\omega_0$ değeri için periyodik\\[1mm]
  Esas frekans $\omega_0$\\[2mm]
  Esas periyot $T_0=\dfrac{2\pi}{\omega_0}$, \ $\omega_0\neq0$
\end{minipage}\hfill
\begin{minipage}[t]{0.53\linewidth}
  \underline{$e^{j\omega_0n}$}\par\smallskip
  $2\pi$'nin katlarındaki $\omega_0$ değerleri için aynı işaret\\[1mm]
  $\dfrac{\omega_0}{2\pi}=\dfrac{m}{N}$ kesirli bir ifade ise periyodik $(m,N>0\in\mathbb{Z})$\\[1mm]
  Esas frekans $\dfrac{\omega_0}{m}$\\[1mm]
  Esas periyot $N_0=m\Bigl(\dfrac{2\pi}{\omega_0}\Bigr)$
\end{minipage}

\begin{ornek}[]
  $x[n]=3e^{j3\pi(n+1/2)/5}$, \quad $y[n]=3e^{j3(n+1/2)/5}$ \quad işaretleri periyodik midir?
\end{ornek}
\begin{cozum}
  \[
    x[n]=3e^{j\frac{3\pi}{5}(n+\frac12)}=3e^{j\frac{3\pi}{10}}\cdot e^{j\frac{3\pi n}{5}}=Ce^{j\omega_0n}
    \ \longrightarrow\
    \begin{array}{l} C=3e^{j3\pi/10}\\[1mm] \omega_0=\dfrac{3\pi}{5}\end{array}
  \]
  \[
    x[n]\overset{?}{=}x[n+N]\ \longrightarrow\ N_0=\Bigl(\frac{2\pi}{\omega_0}\Bigr)m=m\,\frac{2\pi}{\frac{3\pi}{5}}=m\,\frac{10}{3}\in\mathbb{Z}
  \]
  \begin{center} $N_0=10$ periyodu ile periyodik \end{center}
  \[
    y[n]=3e^{j\frac{3}{10}}\cdot e^{j\frac{3n}{5}}=Ce^{j\omega_0n}
    \ \longrightarrow\
    \begin{array}{l} C=3e^{j3/10}\\[1mm] \omega_0=3/5\end{array}
  \]
  \[
    N_0=\frac{2\pi}{3/5}\,m=\frac{10\pi}{3}\,m\notin\mathbb{Z}\ \text{ olup aperiyodik}
  \]
\end{cozum}

\begin{ornek}[]
  $x(t)=2\cos(10t+1)-\sin(4t-1)$ işaretinin periyodu nedir?
\end{ornek}
\begin{cozum}
  \[
    \left.
    \begin{array}{ll}
      A\cos(\omega_0t+\phi)\longrightarrow \omega_{01}=10 & T_{01}=\dfrac{2\pi}{\omega_{01}}=\dfrac{2\pi}{10}=\dfrac{\pi}{5}\\[4mm]
      A\sin(\omega_0t+\phi)\longrightarrow \omega_{02}=4 & T_{02}=\dfrac{2\pi}{\omega_{02}}=\dfrac{2\pi}{4}=\dfrac{\pi}{2}
    \end{array}
    \right\}\ T_0=\pi
  \]
\end{cozum}

\end{document}
